% Nouveau EGA V (cote 25): archive cote 25, site page 60 (PDF page 61).
% typescript. Draft compared: none.
% AI-transcribed by Claude Opus 5.5 (claude-opus-5-5, adaptive thinking, effort high), Message Batches API, batch ega25b (msgbatch_01A7n5YuXubUPNK1zsepi8RG), 2026-11-13; not yet checked by eye.
\scanpage{25/60}
\var{\deleted{\illegible{2}}}{}{deletion}{ega25_25_060-v0001} \mspage{55}

Pour comprendre la signification de ce théorème, il conviendrait
\var{\inserted{\var{\deleted{\illegible{5}}}{}{deletion}{ega25_25_060-v0003}}}{\var{\deleted{\illegible{5}}}{}{deletion}{ega25_25_060-v0003}}{insertion}{ega25_25_060-v0002}
de noter dès le début du nº \hand{(}que  si S=Spec(k)\hand{)}, alors pour une droite
L ``assez générale'' dans $\hand{\check{\ }}P$ ,  la condition de régularité est vérifiée \var{\inserted{(\unsure{notamment pour S.3.|notamment pour 5.3.})}}{(\unsure{notamment pour S.3.|notamment pour 5.3.})}{insertion}{ega25_25_060-v0004}
(cf ra\var{\xout{t}}{}{typewriter strike-out}{ega25_25_060-v0005}trappages signalés au  nº\var{\inserted{\unsure{l'avant}}}{\unsure{l'avant}}{insertion}{ega25_25_060-v0006}\var{\deleted{suivant}}{}{deletion}{ega25_25_060-v0007}).  Par suite, dans les construc%
tions de ``bons'' pinceaux linéaires  indiquées au début du nº, on pouvait
imposer que les pinceaux envisagés \var{\xout{\unsure{étaient}}}{}{typewriter strike-out}{ega25_25_060-v0008} satisfaisaient ladite conditi-
on  (ce qui \var{\xout{ne}}{}{typewriter strike-out}{ega25_25_060-v0009} est une condition\var{\inserted[typed]{de même type mais\hand{,}}}{de même type mais\hand{,}}{insertion}{ega25_25_060-v0010} distincte de celle qui consiste à
exiger que pour tout \hand{$\zeta\in$}L,  $\phi_{\hand{\zeta}}$  soit $X_{k(\hand{\zeta})}$-régulier ). Il fallait inclu
re la condition en question dans les énoncés récapitulatifs préconisés.
\var{\xout{D'autre part, on peut noter}}{}{typewriter strike-out}{ega25_25_060-v0011}  D'autre part, pratiquement  10.2. est
utilisé surtout dans la situation 9.5. ,  ce qui  rend  désirable de
n'énoncer la réformulation de  9.5.  en termes de pinceaux qu'après
10.2. , pour pouvoir  inclure dans l'énoncé en question  également l'iso
morphisme du pinceau avec un éclaté (i.e. donner une descr\var{\xout{u}}{}{typewriter strike-out}{ega25_25_060-v0012}iption de la
situation  permettant une référence commode).  On obtient ainsi un
moyen, pour tout schéma projectif lisse \var{\inserted[typed]{\hand{(}géom \var{\xout{\illegible{1}}}{}{typewriter strike-out}{ega25_25_060-v0014} conn  X  de di\var{\xout{p}}{}{typewriter strike-out}{ega25_25_060-v0015}m$\geq$2 \hand{)}}}{\hand{(}géom \var{\xout{\illegible{1}}}{}{typewriter strike-out}{ega25_25_060-v0014} conn  X  de di\var{\xout{p}}{}{typewriter strike-out}{ega25_25_060-v0015}m$\geq$2 \hand{)}}{insertion}{ega25_25_060-v0013}sur un corps infini k,  de trou\var{\xout{x}}{}{typewriter strike-out}{ega25_25_060-v0016}
ver un\var{\xout{e}}{}{typewriter strike-out}{ega25_25_060-v0017} sous-schéma fermé lisse\var{\xout{ partout}}{}{typewriter strike-out}{ega25_25_060-v0018} \var{\inserted[typed]{non vide}}{non vide}{insertion}{ega25_25_060-v0019} de codimension \var{\deleted{\illegible{1}}}{}{deletion}{ega25_25_060-v0020}2 en chacun de
\var{\xout{puisse être}}{}{typewriter strike-out}{ega25_25_060-v0021}
ses points \var{\xout{(ce qui n'exclut pas qu'il soit vide) et il implique au}}{}{typewriter strike-out}{ega25_25_060-v0022}
\var{\xout{\unsure{contraire si dim X}}}{}{typewriter strike-out}{ega25_25_060-v0023} tel que le schéma éclaté admette une fibration
sur P$^1$ , dont toutes les\hand{/}fibres \var{\inserted[typed]{\hand{(}sont géom intègres, et  dont toutes les fibres \hand{)}}}{\hand{(}sont géom intègres, et  dont toutes les fibres \hand{)}}{insertion}{ega25_25_060-v0024}sont lisses sont au plus un nombre fini,
ces dernières ayant au plus un point \var{\xout{non lisse (et}}{}{typewriter strike-out}{ega25_25_060-v0025} géom. singulier, et
un tel point \var{\xout{est}}{}{typewriter strike-out}{ega25_25_060-v0026} \var{\inserted[typed]{étant}}{étant}{insertion}{ega25_25_060-v0027} rationnel sur k et géom. singulier ordinaire .  Cela
explique l'importance  de l'étude approfondie (\var{\deleted{à peine}}{}{deletion}{ega25_25_060-v0028} \var{\inserted{tout juste}}{tout juste}{insertion}{ega25_25_060-v0029} abordée à l'heure
actuelle) de telles fibrations à fibres singulières, pour réduire dans \var{\xout{x}}{}{typewriter strike-out}{ega25_25_060-v0030}
une certaine mesure l'étude des variétés projectives lisses de dim d
à celle des(familles à un paramètre de) variétés projectives de dim d-1\hand{,}
\hand{pouvant avoir des singularités ordinaires.}

\var{\margin{60}}{\placedmargin{60}}{marginal note placed}{ega25_25_060-v0031}
